% Compute Engine Testing with Privacy-Compliant Production-Like Synthetic Data
% LaTeX reproduction of the VLDB 2024 QDB workshop paper.
% Original: Eric Liu et al., Meta Platforms Inc.
% Proceedings of Workshops at the 50th International Conference on Very Large Data Bases (VLDB), 2024
% Workshop: International Workshop on Quality in Databases (QDB'24)
% Official proceedings PDF: https://vldb.org/workshops/2024/proceedings/QDB/QDB-5.pdf
%
% Compile with: pdflatex synqb.tex (twice for references)

\documentclass[11pt]{article}

\usepackage[margin=1in]{geometry}
\usepackage{amsmath,amssymb}
\usepackage{graphicx}
\usepackage{booktabs}
\usepackage{hyperref}
\usepackage{algorithm}
\usepackage{algpseudocode}

\hypersetup{colorlinks=true,linkcolor=blue,citecolor=blue,urlcolor=blue}

\title{Compute Engine Testing with Privacy-Compliant\\Production-Like Synthetic Data}

\author{
Eric Liu, Jiangnan Cheng, Steve Chuck, Lyublena Antova, Yurgis Baykshtis, Matt David, Ge Gao,\\
Mehrdad Honarkhah, Kuan-Sung Huang, Chen-Kuei Lee, Usman Muhammad, Shihao Peng, Andrii\\
Rosa, Rebecca Schlussel, Michael Shang, Kelvin Silva, Brandon Vo, Zac Wen, Yihao Zhou\\
\small Meta Platforms Inc.\\
\small USA\\
\small \texttt{\{ericyuliu,jncheng\}@meta.com}
}

\date{}

\begin{document}

\maketitle

\begin{abstract}
Compute engines require thorough and continuous testing where
synthetic benchmarks or hand-crafted queries and data sets are not
enough, yet accessing customer data and queries either is impossible
or brings a big burden of user data protection. At Meta, we build
SynQB - a query bank for compute engine testing that leverages
privacy-compliant production-like synthetic data with good
data quality, generated by a carefully designed differentially-private
synthetic data generation (DPSDG) algorithm. This solution solves
our common pain points in testing by enabling test coverages for
different compute engines, query operators and workloads and facilitating
various correctness and performance regression detections.
As a result, SynQB greatly improves our compute engine release
confidence and enhances the reliability of our data warehouse.
\end{abstract}

\noindent\textbf{PVLDB Reference Format:}\\
Eric Liu, Jiangnan Cheng, Steve Chuck, Lyublena Antova, Yurgis Baykshtis,
Matt David, Ge Gao, Mehrdad Honarkhah, Kuan-Sung Huang, Chen-Kuei
Lee, Usman Muhammad, Shihao Peng, Andrii Rosa, Rebecca Schlussel,
Michael Shang, Kelvin Silva, Brandon Vo, Zac Wen, Yihao Zhou. Compute
Engine Testing with Privacy-Compliant Production-Like Synthetic Data.
\textit{PVLDB}, 14(1): XXX-XXX, 2020.\\
doi:XX.XXX/XXX.XX

\section{Introduction}
Compute engines are critical in data-driven organizations, especially
when large amounts of data need to be processed to generate
insights and support decision-making. As such, they require thorough
and continuous testing to ensure correctness, performance
and reliability. As in any other domain, simulating realistic workloads
during testing helps with the proactive discovery of bugs and
inefficiencies before they hit the database. A big challenge when
testing compute engines is that database vendors often do not have
access to customer data and queries, or are not allowed to make
copies of these outside of the customers' environment. To address
this challenge, vendors have historically relied on one or both of
the following: synthetic benchmarks or hand-crafted queries and
data sets to emulate a customer environment. There is a variety
of synthetic query and data generators attempting to create a true
representation of real customer workloads, including TPC-H and
TPC-DS among others. Those however are only a start, as they
fail to represent realistic customer workloads. Another approach
taken by database vendors is to sample customer queries and workloads
and use these for testing. While this provides a much more
realistic test bed and removes many of the limitations of testing
with synthetic benchmarks, it has limited applicability as it requires
customers to approve using their data and queries in a test environment,
and puts a big burden on the database vendor to ensure
customer data is protected. At Meta, preserving user data privacy is
critical and tables are subject to strict restrictions on user access, as
well as data retention policies. This presents a challenge to testing
with real workloads.

We built SynQB to solve the challenge of testing realistic workloads
in a privacy compliant way. SynQB is a query bank framework
for compute engine testing that leverages synthetic data generation
(SDG) that provides production-like synthetic data for testing.
SDG is an emerging technique that generates simulated data with
similar statistical properties and characteristics of the real data
while protecting the sensitive information of each individual user.
It was noted by the National Science and Technology Council a ``key
technical approach for privacy-preserving data sharing and analytics''
[1]. In particular, our designed SDG algorithm is protected by
differential privacy (DP) [5, 6], a state-of-the-art privacy standard
that guarantees the identification of a specific user is technically
impossible. The production-like and privacy-compliant synthetic
data are of high data quality, solving our common pain points in
testing by enabling test coverage for different compute engines,
query operators and workloads and facilitating detection of correctness
and performance regressions. As a result, SynQB greatly
improves our compute engine release confidence and enhances the
reliability of our data warehouse.

In light of prior work, our contributions are as follows. First, we
design and implement the system architecture of SynQB, which
includes a creation workflow where production queries are rewritten,
production-like privacy-compliant synthetic tables are generated,
and a usage suite is generated which incorporates various
correctness and performance regression detections. Second, we
design a generic differentially-private synthetic data generation
(DPSDG) algorithm for compute engine testing, which handles
user-identifiable information (UII) and user features differently,
and can also be parallelized across multiple machines and CPUs. Third,
through evaluations we validate that SynQB generates high-quality
synthetic data that provide good signals for compute engine testing.

\section{Related Work}

\subsection{Compute Engines at Meta}
Presto [15] is a distributed SQL query engine used for low-latency
interactive workloads as well as long-running ETL pipelines involving
exabyte-scale data volumes at Meta. It is widely used by teams
and systems across the company for diverse analytical use cases
such as business intelligence, machine learning, feature engineering,
experimentation platforms etc.

In recent years, Meta has successfully implemented several Presto
evolutions like Presto-on-Velox and Presto-on-Spark, which can be
seen as extensions of the classic Presto engine. These engines integrate
Presto with other big data processing technologies: Apache
Spark in the case of Presto-on-Spark, and Velox vectorized execution
engine [12] for Presto-on-Velox to achieve enhanced query
performance, reliability and scalability for specific data analytics
workloads. One noteworthy aspect is that these engines all share
the same Presto SQL dialect syntax and semantics. The unified language
interface, in addition to offering a consistent user experience
and reducing user friction, has helped improve the development
of the snapshot testing framework. This is achieved by promoting
query bank reusability and seamless cross-engine verification as
queries designed for Presto can easily be adapted to run on other
engines sharing the same interface.

\subsection{Compute Engine Testing}
Presto Verifier\footnote{\url{https://prestodb.io/docs/current/admin/verifier.html}}
is an open-source testing tool used for executing
queries and ensuring their correctness by comparing the results
with a control run. Typically, the control side represents the stable
production environment while the test side represents the release
candidate environment. This process involves checking the consistency
of the checksums generated by both query results: any
discrepancies with deterministic queries indicate correctness regressions.
However, a potential drawback of utilizing Presto Verifier
is the production data source may change overtime or even get
removed entirely, resulting in unreliable and low-signal test results.

SnowTrail [18] is Snowflake's testing framework that leverages
its time travel capability to execute queries at predefined timestamps,
thereby ensuring data consistency. However, this framework
faces challenges due to constantly evolving schema or data changes
that may render original queries ineffective overtime.

SynQB addresses this concern by employing reliable, production-like,
and privacy-compliant synthetic data, thereby minimizing data
variability and providing a more robust foundation for effective
and high signal testing, encompassing both correctness and performance
regression detections. This method offers a cleaner solution
compared to the previous approaches. Besides, SynQB offers advantages
by not negatively affecting the customers' production traffic
in two key aspects: 1) By minimizing the necessity for control side
runs on the production cluster; 2) By limiting the requirement for
both control and test sides to utilize customer production data.

\subsection{Synthetic Benchmarks}
There is a long history of synthetic benchmarks for database testing
with a large focus on performance testing. The Transaction Processing
Council (TPC)\footnote{\url{https://www.tpc.org/}} maintains a number of benchmarks for various
domains, including TPC-C, TPC-H, TPC-DS, and the more recent
TPCx-AI [4]. In the data analytics space TPC-H and TPC-DS are the
main benchmarks. They collectively feature a little over 120 queries
and provide a data generator to produce sample data for testing.
While these are often a good start, synthetic generators fall short
in modeling true customer workloads. For example, neither TPC-H
or TPC-DS support complex data types such as arrays or maps, or
test for operations on these. In addition, the research community
has contributed a number of proposals for synthetic generators,
including schema- and query-aware generators [7, 10, 13, 14], but
these proposals do not necessarily have privacy in mind.

\subsection{DPSDG}
There are a wide variety of works designing SDG algorithms with
good utility under DP guarantee. Based on the underlying model,
they can be split into two main categories: 1) \textbf{Marginal-based
models} (such as AIM [11] and DP Gaussian Copula Kendall [9])
measure the marginal distributions of the input dataset with DP
guarantee, and then generate synthetic data that comply with the
marginal distributions; 2) \textbf{GAN-based models} (such as DPGAN
[17] and PATE-GAN [8]) use differentially-private stochastic gradient
descent (DPSGD) method [2] to train a generative adversarial
net (GAN), which consists of a generator for synthetic data. In most
of these works, the utility metrics of synthetic data were defined
as marginal distribution, pairwise correlation and the accuracy of
machine learning inference of the synthetic data, etc. To the best of
our knowledge, this paper is the first work to consider the utility
of synthetic data in the context of compute engine testing.

\section{System Architecture}
In this chapter we first introduce the creation workflow of SynQB
in Sec.~3.1 and then the corresponding usage suite in Sec.~3.2. The
advantages of SynQB are highlighted in Sec.~3.3.

\subsection{SynQB Creation Workflow}

\begin{figure}[h]
\centering
\fbox{\parbox{0.9\linewidth}{\centering\textit{[Figure 1: SynQB creation workflow diagram]}\\
\small Prod Query Pool $\to$ Query Selector $\to$ Query Parser
$\to$ (Prod Tables $\to$ DPSDG Fit/Gen $\to$ Synthetic Tables;
Query Syntax Tree $\to$ Query Rewriter $\to$ Rewritten Query)
$\to$ SynQB $\to$ Persist result checksum as Baseline}}
\caption{SynQB creation workflow. Queries are selected from
production workload and rewritten with generated privacy-compliant
and production-like tables.}
\label{fig:creation}
\end{figure}

The SynQB creation workflow, as illustrated in Fig.~\ref{fig:creation},
encompasses the following three steps.

\textbf{Query selection.} The source of SynQB is a production query
pool. We adopt a query selector with various selection criteria based
on the coverage requirements. There are two types of coverage we
consider: 1) \textbf{Generic code coverage}, which further encompasses
several major dimensions, such as coverage of different Presto operators
including TableScan, Join, Aggregation, TableWriter, etc. and
coverage of different data structures including map, array, struct,
JSON, etc.; 2) \textbf{Important workloads coverage}, such as company-wide
critical dashboards, ads experiments workloads, etc.

\textbf{Query rewriting and synthetic table generation.} A selected
query will then be parsed by a query parser, which extracts the production
table names and produces the corresponding query syntax
tree. For each production table, we leverage a DPSDG algorithm to
generate a production-like and privacy-compliant synthetic table.
The algorithm takes two steps: 1) a \textit{fit} step that trains a compressed
model $M$ from the input production table $D$ with $(\epsilon,\delta)$-DP
(defined in Sec.~4.1) guarantee, denoted as
\begin{equation}
M = \text{DPSDG-Fit}(D; \epsilon, \delta),
\end{equation}
and 2) a \textit{generate} step that produces a synthetic table $D_{\text{syn}}$
with $N_{\text{syn}}$ number of rows according to the model $M$, denoted as
\begin{equation}
D_{\text{syn}} = \text{DPSDG-Gen}(M, N_{\text{syn}}).
\end{equation}
The algorithm behind these two steps will be discussed in details
in Sec.~4.2. Such two-step process enables parallelized generation
through vertical scaling and horizontal scaling. After synthetic table
generation, the query will also be rewritten with the original query
syntax tree and the new synthetic table names.

\textbf{Query and synthetic table persistence.} The rewritten query
and the synthetic tables will be added to the SynQB. The query will
also be run with a stable version of Presto, and then the checksum
of the output will be persisted as the baseline, which can be used
in the future test runs for correctness regression detection.

\subsection{SynQB Usage Suite}

\begin{figure}[h]
\centering
\fbox{\parbox{0.9\linewidth}{\centering\textit{[Figure 2: SynQB usage suite diagram]}\\
\small SynQB Usages $\to$ Correctness Regression Detection
(Detect bugs $\to$ Release Verification, Major upgrades, Cross Engine Verification);
Performance Regression Detection (Release Verification $\to$
Continuous runs to generate perf baseline $\to$ Detect regressions by
comparing w/ baseline; New Optimizations $\to$ A/B Testing, Verify
improvements, Avoid regression)}}
\caption{SynQB usage suite. SynQB can be used to detect
various kinds of regressions.}
\label{fig:usage}
\end{figure}

As shown in Fig.~\ref{fig:usage}, SynQB is mainly used to detect regressions,
including correctness and performance regressions.

\textbf{Correctness regression detection.} For compute engines, the
most important requirement is to guarantee that results are correct.
Presto, Presto-On-Velox, and Presto-On-Spark all share the same
SQL interface, meaning that the same query is expected to work
on these engines interchangeably. Therefore, we use the same test
query to not only verify continuous releases for each compute
engine itself, but also cross check each other. For example, while
developing Presto-on-Velox, we compare its query results with
baseline result generated by stable Java-based Presto, which has
not only exposed correctness issues but also has helped identify
missing functionality in the new engine.

\textbf{Performance regression detection.} Queries and tables created
for correctness regression detection can also be used for performance
regression detection. By using immutable synthetic data
and a consistent testing environment, we gather performance results
that provide stronger signals than traditional shadow testing
solutions and are more representative of our workload than industry
benchmarks. The key strategy here is to establish a consistent
and trustworthy baseline, which facilitates effective and meaningful
regression detection analysis. To accomplish this, we run tests
multiple times per day across releases, recording key metrics for
each query, such as CPU usage, memory usage, and execution time.
Based on this high signal setup, for major upgrades and optimization,
we also conduct A/B testing to measure the improvements and
avoid accidental regressions. This approach allows us to compare
the performance of different versions under controlled conditions,
ensuring that enhancements deliver the expected benefits.

\subsection{Highlights}
We highlight some advantages of our SynQB as follows.

\textbf{Sustainable testing with synthetic data.} To ensure production-like
quality and immutability for sustainable testing purposes, we
generate privacy-compliant synthetic data. SynQB transforms production
queries to use these immutable test datasets as inputs. This
approach addresses the issue of unreliable data sources from the
ground up, enabling us to develop a sustainable and high-signal
test framework.

\textbf{Early signals.} From daily testing results, we have proactively
prevented correctness regressions from slipping into production,
effectively averting potential data corruptions. Furthermore, SynQB
offers expedited feedback compared to conventional production
shadow testing solutions. This efficiency is achieved through the
manipulation of synthetic test data sizes. We use smaller data sets
for correctness testing to decreases the runtime of end-to-end testing.
Additionally, the short execution time enables the integration
of SynQB into the pull-request landing process. This enhancement
can accelerate the testing cycle and facilitate the early detection and
resolution of bugs. Developers can receive immediate feedback on
their changes, promoting a more efficient and iterative development
process.

\textbf{Reliable performance regression A/B testing framework.}
The use of synthetic data fundamentally reduces variability at the
data layer, thereby providing a more robust foundation for reliable
testing. Our performance regression detection framework is
particularly relevant for projects that involve significant changes
or upgrades, such as the introduction of new file formats in data
warehouses, SQL query optimizations, and major library upgrades.
By using this approach, we ensure that each enhancement not only
meets the intended functionality improvements but also maintains
or enhances overall system performance.

\textbf{Historical failed queries.} We systematically collect queries
that had production regressions and incorporate them into a historical-failed
queries test suite. This suite is then used for verifying future
releases, ensuring that previously encountered issues are adequately
addressed and do not recur.

\textbf{Flexible and pluggable testing framework.} SynQB is highly
extensible, facilitating integration with each compute engine's test
runner. Presently, it supports a range of engines including Presto,
Presto-on-Velox, and Spark. Beyond compute engines, SynQB's
synthetic data can also be directly utilized for flexible integration
testing and library testing across different platforms.

\textbf{Minimized Regional Capacity Limitations.} In Meta's data
warehouse, queries are run on servers geographically near to where
the data is stored to ensure low latency and network IO. As such,
shadow testing solutions that run on real customer data require test
environments in all regions where customer data may be located,
which is not always feasible. In contrast, SynQB employs synthetic
data that is not bound by regional constraints since the data can
be automatically replicated to any region. This flexibility makes
it possible to leverage any available test cluster for testing any
workload.

\section{DPSDG Deep Dive}
In this chapter, we first briefly introduce some basic definitions and
theorems of DP in Sec.~4.1, and then give an in-depth explanation
of our DPSDG algorithm for SynQB in Sec.~4.2.

\subsection{DP Basics}
Consider a dataset $D$ with $N$ records (i.e., rows for a tabular dataset).
We use $\mathcal{D}$ to denote the domain of $D$.

Next we formally define DP, which protects user-level privacy
based on the notation of neighboring datasets, i.e., a pair of datasets
which differ in just one record.

\textit{Definition 4.1 (Differential Privacy (DP) [6]).} A randomized algorithm
$\mathcal{M} : \mathcal{D} \to \mathcal{R}$ is said to satisfy $(\epsilon,\delta)$-DP if
$\forall D, D' \in \mathcal{D}$ s.t.\ they are neighboring datasets (denoted by
$D \sim D'$), and any possible outcome $\mathcal{S} \subseteq \mathcal{R}$, we always have
\begin{equation}
\Pr(\mathcal{M}(D) \in \mathcal{S}) \leq e^{\epsilon}\Pr(\mathcal{M}(D') \in \mathcal{S}) + \delta.
\end{equation}

Essentially, with DP guarantee, one cannot readily differentiate
whether the input of algorithm $\mathcal{M}$ is $D$ or $D'$, due to the randomness
of the algorithm. The combination $(\epsilon, \delta)$ is referred to as privacy
budget, which measures the privacy level of the algorithm $\mathcal{M}$.

Next, we introduce subsample procedure for dataset $D$ and what
it implies under the context of DP.

\textit{Definition 4.2 (Subsample Procedure [16]).} For dataset $D$ with $N$
records, the \texttt{subsample} procedure selects a random dataset $\tilde{D}$ from
the uniform distribution over all subsets of $D$ of size $\tilde{N}$. The ratio
$\gamma := \tilde{N}/N$ is defined as the sampling parameter of the \texttt{subsample}
procedure.

\textit{Theorem 4.3 (Differential Privacy with subsampling [16]).}
\textit{If $\mathcal{M}$ is $(\epsilon,\delta)$-DP, then
$\mathcal{M} \circ \texttt{subsample}$ is $(\log(1+\gamma(e^{\epsilon}-1)), \gamma\delta)$-DP.}

According to the above theorem, if we target at $(\epsilon,\delta)$-DP, then
we can apply an algorithm $\tilde{\mathcal{M}}$ with only
$(\log(1+\frac{1}{\gamma}(e^{\epsilon}-1)), \delta/\gamma)$-DP
to the subsampled dataset $\tilde{D}$ as the input.

Below we introduce two widely-used properties of DP.

\textit{Theorem 4.4 (Sequential Composition [6]).}
\textit{If randomized algorithm $\mathcal{M}_i : \mathcal{D} \to \mathcal{R}_i$
satisfies $(\epsilon_i,\delta_i)$-DP, $\forall i \in \{1,2,\cdots,k\}$, then
randomized algorithm $\mathcal{M}(D) \triangleq
(\mathcal{M}_1(D),\mathcal{M}_2(D),\cdots,\mathcal{M}_k(D))$
satisfies $(\sum_i \epsilon_i, \sum_i \delta_i)$-DP.}

\textit{Theorem 4.5 (Post Processing [6]).}
\textit{For randomized algorithm $\mathcal{M} : \mathcal{D} \to \mathcal{R}$
and algorithm $g : \mathcal{R} \to \mathcal{G}$ (either randomized or
deterministic), if $\mathcal{M}$ is $(\epsilon,\delta)$-DP, then
$g \circ \mathcal{M}$ is also $(\epsilon,\delta)$-DP.}

Lastly, we introduce a variant of Gaussian mechanism which protects
a deterministic function with DP by adding Gaussian noises.

\textit{Theorem 4.6 (Analytic Gaussian Mechanism [3]).}
\textit{Analytic Gaussian mechanism is a $(\epsilon,\delta)$-DP mechanism
that adds Gaussian noise to the result of a sensitive deterministic
function $g : \mathcal{D} \to \mathbb{R}^p$:}
\begin{equation}
\mathcal{M}(D) = g(D) + \sigma\mathcal{N}(0, \mathbb{I}_p)
\end{equation}
\textit{where $\sigma$ is the scale of the Gaussian noise that satisfies}
\begin{equation}
\Phi\!\left(\frac{\Delta_g}{2\sigma} - \frac{\epsilon\sigma}{\Delta_g}\right)
- e^{\epsilon}\Phi\!\left(-\frac{\Delta_g}{2\sigma} - \frac{\epsilon\sigma}{\Delta_g}\right) \leq \delta
\end{equation}
\textit{where $\Phi$ is the CDF of the standard univariate Gaussian distribution
and $\Delta_g \triangleq \sup_{D \sim D'} \|g(D) - g(D')\|$ is the $L_2$-sensitivity of $g$.}

\subsection{DPSDG Algorithm}
We notice that a user dataset generally contain two categories of information:
1) \textbf{User features}, i.e., the demographics of an individual,
such as country, gender, age, etc; and 2) \textbf{UII}, i.e., the information
that can be used to directly identify a specific user, such as user ID,
phone number, name, etc.

An implicit assumption of most prior works is that releasing the
\textit{values} of the columns (e.g., users are from these countries: USA and
Canada) is not risky, but the \textit{distributions} of the columns (e.g., 70\%
and 30\% of the users are from USA and Canada, respectively) is
sensitive and hence should be protected by DP. For feature columns,
such assumption normally makes sense, because it is quite common
that a non-trivial number of people have the same demographics.
However, these prior works didn't consider UII, whose \textit{values} (e.g.,
there exists a user with ID 1001) are inherently very sensitive, and
thus should be handled differently.

Therefore, in this paper, as illustrated in Table~\ref{tab:split}, we split $D$
vertically: $D = (D_{\text{ft}}, D_{\text{uii}})$, where $D_{\text{ft}}$ and $D_{\text{uii}}$
are datasets with feature columns and UII columns, respectively. In the following
sections, we will apply different DPSDG algorithms to $D_{\text{ft}}$ and $D_{\text{uii}}$.
In particular, we carefully designed a DPSDG algorithm for UII such
that we can preserve the distribution of UIIs under DP, while the
corresponding values will not be leaked.

\begin{table}[h]
\centering
\caption{Split of Dataset $D$}
\label{tab:split}
\begin{tabular}{ccccc}
\toprule
\multicolumn{2}{c}{Feature Columns ($D_{\text{ft}}$)} &&
\multicolumn{2}{c}{UII Columns ($D_{\text{uii}}$)} \\
\cmidrule{1-2}\cmidrule{4-5}
Country & Gender & $\cdots$ & User ID & Name & $\cdots$ \\
\midrule
USA & Male && 1001 & John \\
Canada & Female && 1002 & Alice \\
\bottomrule
\end{tabular}
\end{table}

\textbf{DPSDG Algorithm for feature columns.} Our DPSDG algorithm
for feature columns is presented in Algorithm~\ref{alg:feature}, which is
largely based on the DP Gaussian Copula Kendall from [9], with
some modifications. For the fit step, in line 2, we first flatten the
structural columns and discretize the continuous columns, which
is because DP Gaussian Copula Kendall, as a marginal-based mechanism,
requires the input to be unnested and discretized in order
to compute the discrete marginal distribution; Next, in line 3-6, we
split the privacy budget and then compute the noisy empirical 1-d
marginal distribution $F_{\text{ft}}$ and Pearson's correlation coefficient matrix
$\Sigma_{\text{ft}}$ (estimated from Kendall's $\tau$ correlation coefficient matrix
$\Sigma^{\tau}_{\text{ft}}$, see details in [9]); and in line 7, the trained model includes
these noisy marginals and will be returned. For the generate step,
we sample $N_{\text{syn}}$ samples that comply with these noisy marginals
in line 9, and then inverse transform the discretized continuous
columns and then flattened structural columns in line 10, which
produces the synthetic feature dataset $D_{\text{ft},\text{syn}}$. This algorithm can
be proved to be $(\epsilon,\delta)$-DP based on Theorem 4.4 and Theorem 4.5.

\begin{algorithm}[h]
\caption{Differentially-Private Synthetic Data Generation for feature columns DPGCKendall1}
\label{alg:feature}
\begin{algorithmic}[1]
\Procedure{DPGCKendall1-Fit}{$D_{\text{ft}}, \epsilon, \delta$}
\State Flatten the structural columns and then discretize the continuous columns
\State Split the privacy budget into $\epsilon_1, \delta_1$ and $\epsilon_2, \delta_2$,
  s.t.\ $\epsilon_1 + \epsilon_2 = \epsilon$, and $\delta_1 + \delta_2 = \delta$
\State Compute the empirical 1-d marginal distribution $F_{\text{ft}}$ and
  Kendall's $\tau$ correlation coefficient matrix $\Sigma^{\tau}_{\text{ft}}$ of $D_{\text{ft}}$
\State Apply analytic Gaussian mechanism to $F_{\text{ft}}$ and $\Sigma^{\tau}_{\text{ft}}$,
  making them $(\epsilon_1,\delta_1)$-DP and $(\epsilon_2,\delta_2)$-DP, respectively
\State Estimate the Pearson's correlation coefficient matrix $\Sigma_{\text{ft}}$
\State \Return $M_{\text{ft}} = (F_{\text{ft}}, \Sigma_{\text{ft}})$
\EndProcedure
\Procedure{DPGCKendall1-Gen}{$M_{\text{ft}}, N_{\text{syn}}$}
\State Sample $D_{\text{ft},\text{syn}}$ with $N_{\text{syn}}$ samples that comply with $F_{\text{ft}}, \Sigma_{\text{ft}}$
\State Inverse transform the discretized continuous columns and then
  flattened structural columns of $D_{\text{ft},\text{syn}}$
\State \Return $D_{\text{ft},\text{syn}}$
\EndProcedure
\end{algorithmic}
\end{algorithm}

\textbf{DPSDG Algorithm for UII columns.} Our DPSDG algorithm
for UII columns is shown in Algorithm~\ref{alg:uii}. The key of the algorithm
is to compute the second-order histogram of $D_{\text{uii}}$, denoted by $H_{\text{uii}}$,
as in line 2. For example, consider a UII dataset $D_{\text{uii}}$ with just one
column ``User ID''. As illustrated in Fig.~\ref{fig:hist}, we first compute the
first-order histogram of $D_{\text{uii}}$, which indicates the frequency (or
occurrence) of each user ID. Due to the sensitivity of the values
of user IDs themselves, we should not directly preserve the first-order
histogram. Instead, we further compute the second-order
histogram (i.e., the histogram of the first-order histogram) of $D_{\text{uii}}$,
which contains only the values of the frequencies as well as the
frequencies of these frequencies. This way, we not only drop those
sensitive ID values, but also compress the key information since
the cardinality of user IDs is in general much larger than their
frequencies. Once we have $H_{\text{uii}}$, the remaining part of the algorithm
(line 3-7, i.e., noise perturbation and sampling) is quite similar to
Algorithm~\ref{alg:feature}, except that we need to sample $D_{\text{uii},\text{syn}}$ from a test UII
dataset domain $\mathcal{D}^{\text{test}}_{\text{uii}}$, i.e., synthetic UIIs are essentially pseudonyms
unlinkable to the original UIIs. In the example of Fig.~\ref{fig:hist}, the synthetic
user IDs start from 2001, and cannot be linked back to the real user
IDs in the input dataset.

\begin{algorithm}[h]
\caption{Differentially-Private Synthetic Data Generation for UII columns DPUII}
\label{alg:uii}
\begin{algorithmic}[1]
\Procedure{DPUII-Fit}{$D_{\text{uii}}, \epsilon, \delta$}
\State Compute the second-order histogram $H_{\text{uii}}$ of $D_{\text{uii}}$
\State Apply analytic Gaussian mechanism (Theorem 4.6) to $H_{\text{uii}}$,
  making it $(\epsilon,\delta)$-DP
\State \Return $M_{\text{uii}} = H_{\text{uii}}$
\EndProcedure
\Procedure{DPUII-Gen}{$M_{\text{uii}}, N_{\text{syn}}$}
\State Sample $D_{\text{uii},\text{syn}}$ from a test domain $\mathcal{D}^{\text{test}}_{\text{uii}}$
  with $N_{\text{syn}}$ samples that comply with $H_{\text{uii}}$
\State \Return $D_{\text{uii},\text{syn}}$
\EndProcedure
\end{algorithmic}
\end{algorithm}

\begin{figure}[h]
\centering
\fbox{\parbox{0.9\linewidth}{\centering\textit{[Figure 3: Illustration of second-order histogram computation]}\\
\small First-Order Histogram (Production): User ID $\to$ Freq.\
(1001: 4, 1002: 2, 1003: 2, 1004: 1, 1005: 1) $\to$
Second-Order Histogram (Production): Freq.\ $\to$ Freq.\ of Freq.\
(1: 2, 2: 2, 3: 0, 4: 1) $\to$
Second-Order Histogram (Synthetic): (1: 1, 2: 2, 3: 1, 4: 1) $\to$
First-Order Histogram (Synthetic): User ID $\to$ Freq.\
(2001: 1, 2002: 4, 2003: 3, 2004: 2, 2005: 2)}}
\caption{Illustration of second-order histogram computation. The left two tables are the first-order and second-order histograms for production table, and the right two tables are second-order and first-order histograms for synthetic table.}
\label{fig:hist}
\end{figure}

\textbf{Complete DPSDG Algorithm.} The complete DPSDG algorithm
to synthesize a dataset $D$ is given in Algorithm~\ref{alg:complete}, which incorporates
privacy budget splitting, subsampling and utilizing Algorithm~\ref{alg:feature}
and~\ref{alg:uii}. According to Theorem 4.3, 4.4 and 4.5, the algorithm can
be proved to be $(\epsilon,\delta)$-DP.

\begin{algorithm}[h]
\caption{Differentially-Private Synthetic Data Generation DPSDG}
\label{alg:complete}
\begin{algorithmic}[1]
\Procedure{DPSDG-Fit}{$D, \epsilon, \delta$}
\State Split the privacy budget into $\epsilon_{\text{ft}}, \delta_{\text{ft}}$ and
  $\epsilon_{\text{uii}}, \delta_{\text{uii}}$, s.t.\ $\epsilon_{\text{ft}} + \epsilon_{\text{uii}} = \epsilon$,
  and $\delta_{\text{ft}} + \delta_{\text{uii}} = \delta$
\State Subsample $D_{\text{ft}}$ and $D_{\text{uii}}$ with rate $\gamma_{\text{ft}}$ and $\gamma_{\text{uii}}$
  respectively. Get subsampled dataset $\tilde{D}_{\text{ft}}$ and $\tilde{D}_{\text{uii}}$
\State Adjust $\epsilon_{\text{ft}}, \delta_{\text{ft}}$ and $\epsilon_{\text{uii}}, \delta_{\text{uii}}$
  according to Theorem 4.3. Get $\tilde{\epsilon}_{\text{ft}}, \tilde{\delta}_{\text{ft}}$
  and $\tilde{\epsilon}_{\text{uii}}, \tilde{\delta}_{\text{uii}}$
\State $M_{\text{ft}} = \text{DPGCKendall1-Fit}(\tilde{D}_{\text{ft}}; \tilde{\epsilon}_{\text{ft}}, \tilde{\delta}_{\text{ft}})$
\State $M_{\text{uii}} = \text{DPUII-Fit}(\tilde{D}_{\text{uii}}; \tilde{\epsilon}_{\text{uii}}, \tilde{\delta}_{\text{uii}})$
\State \Return $M = (M_{\text{ft}}, M_{\text{uii}})$
\EndProcedure
\Procedure{DPSDG-Gen}{$M, N_{\text{syn}}$}
\State $D_{\text{ft},\text{syn}} = \text{DPGCKendall1-Gen}(M_{\text{ft}}, N_{\text{syn}})$
\State $D_{\text{uii},\text{syn}} = \text{DPUII-Gen}(M_{\text{uii}}, N_{\text{syn}})$
\State \Return $D_{\text{syn}} = (D_{\text{ft},\text{syn}}, D_{\text{uii},\text{syn}})$
\EndProcedure
\end{algorithmic}
\end{algorithm}

\section{Evaluation}
In this section, we evaluate the quality of the synthetic data by comparing
it with production data, with a focus on the following utility metrics
that are relevant in the context of compute engine testing.

\textbf{Data volume.} Fig.~\ref{fig:eval} shows the number of rows in the synthetic
data generated with different sample sizes. With our design, the
sample size can be adjusted to accommodate testing requirements.
For instance, small datasets allow for faster testing to provide
early signals during development, while large datasets can be used
for thorough verification.

\begin{figure}[h]
\centering
\fbox{\parbox{0.9\linewidth}{\centering\textit{[Figure 4: Synthetic data volume evaluation]}\\
\small Table1: production vs synthetic row counts at
$10^4, 10^5, 10^6, 10^7$; Table2: production vs synthetic at
$10^5, 10^6, 10^7$ (log-scale bar charts showing synthetic
volumes matching production volumes)}}
\caption{Number of rows in the synthetic data generated with different sample sizes.}
\label{fig:eval}
\end{figure}

\textbf{Data distribution.} Due to the two-step fit/generate process
of our DPSDG algorithm (introduced in Sec.~3.1), we are able to
parallelize the generation across multiple machines and CPUs. Also,
we are able to sample synthetic data with the same number of rows
as the input production data, and then check whether the production-like
distribution is captured by comparing the actual data with
expected data, which is generated by a random sampler. Furthermore,
data distribution has a direct impact on query performance, making
it an important aspect for performance regression detection.

\textbf{Query checksum.} We run the rewritten test queries against both
the synthetic and production data using the stable version of Presto,
and verify that the results match (or are similar) as measured by
checksums, which confirms the correctness of the synthetic data
generation process.

\section{Conclusions}
In this paper, we designed and built SynQB, a query bank framework
for compute engine testing that leverages synthetic data generation.
SynQB provides production-like and privacy-compliant synthetic
data that solves our common pain points in testing by enabling
test coverages for different compute engines, query operators and
workloads and facilitating detection of correctness and performance
regressions. In particular, we designed a generic differentially-private
synthetic data generation (DPSDG) algorithm for SynQB,
which can be proved to satisfy $(\epsilon,\delta)$-DP and can handle
user-identifiable information and user features differently. Finally,
our evaluation results show that the synthetic data generated by
SynQB is of high data quality that provides good signals for
compute engine testing.

\begin{thebibliography}{18}

\bibitem{ref1}
National Science and Technology Council. National strategy
to advance privacy-preserving data sharing and analytics.
\url{https://www.whitehouse.gov/wp-content/uploads/2023/03/National-Strategy-to-Advance-Privacy-Preserving-Data-Sharing-and-Analytics.pdf}, 2023.

\bibitem{ref2}
M.~Abadi, A.~Chu, I.~Goodfellow, H.~B.~McMahan, I.~Mironov,
K.~Talwar, and L.~Zhang. Deep learning with differential privacy.
In \textit{Proceedings of the 2016 ACM SIGSAC Conference on Computer and
Communications Security}, CCS '16, page 308--318, New York, NY, USA, 2016.
Association for Computing Machinery.

\bibitem{ref3}
B.~Balle and Y.-X.~Wang. Improving the gaussian mechanism for
differential privacy: Analytical calibration and optimal denoising.
In \textit{International Conference on Machine Learning}, pages 394--403.
PMLR, 2018.

\bibitem{ref4}
R.~Nambiar and M.~Poess. The making of tpc-ds. In \textit{Proceedings of the
32nd International Conference on Very Large Data Bases}, VLDB '06,
page 1049--1058. VLDB Endowment, 2006.

\bibitem{ref5}
C.~Dwork. Differential privacy. In \textit{International Colloquium on
Automata, Languages, and Programming}, pages 1--12. Springer, 2006.

\bibitem{ref6}
C.~Dwork, A.~Roth, et~al. The algorithmic foundations of differential
privacy. \textit{Foundations and Trends in Theoretical Computer Science},
9(3-4):211--407, 2014.

\bibitem{ref7}
M.~Hao, X.~Li, X.~Meng, and S.~Ma. Efficient and effective
genetic algorithms for covering array generation. \textit{IEEE Transactions
on Software Engineering}, 42(9):873--895, 2016.

\bibitem{ref8}
J.~Jordon, J.~Yoon, and M.~Van~Der~Schaar. Pate-gan: Generating
synthetic data with differential privacy guarantees. In \textit{International
Conference on Learning Representations}, 2018.

\bibitem{ref9}
H.~Li, L.~Xiong, and X.~Jiang. Differentially private synthesization
of multi-dimensional data using copula functions. In \textit{Advances in
Database Technology-EDBT 2014: 17th International Conference on
Extending Database Technology, Athens, Greece, March 24-28, 2014.},
pages 475--486, 2014.

\bibitem{ref10}
P.~Loiseaux, L.~Berti-Equille, et~al. Cost-effective bulk loading
with sql server integration services and ibm db2. In \textit{ICDE},
pages 1457--1458, 2007.

\bibitem{ref11}
R.~McKenna, D.~Sheldon, and G.~Miklau. Graphical-model based
estimation and inference for differential privacy. In \textit{International
Conference on Machine Learning}, pages 4435--4444. PMLR, 2019.

\bibitem{ref12}
P.~Pedreira, O.~Bhardwaj, A.~Rai, M.~Dey, K.~Saha, G.~Chen,
S.~Wang, L.~Xia, M.~Aston, A.~Chen, et~al. Velox: Meta's unified
execution engine. \textit{Proceedings of the VLDB Endowment},
15(12):3372--3384, 2022.

\bibitem{ref13}
F.~Russo dos Santos, P.~Sottile, and E.~Muller. An effective and
efficient approach for big data forensics. In \textit{EDBT}, pages 584--595,
2021.

\bibitem{ref14}
T.~Rabl, F.~Llirbat, et~al. Data generator for the social network
benchmark. 2017.

\bibitem{ref15}
R.~Sethi, M.~Traverso, D.~Sundstrom, D.~Phillips, W.~Xie, Y.~Sun,
N.~Yegitburu, D.~Feldman, H.~Shi, and C.~Wang. Presto: Sql on
everything. In \textit{2019 IEEE 35th International Conference on Data
Engineering (ICDE)}, pages 1802--1813. IEEE, 2019.

\bibitem{ref16}
N.~Srinivasa, B.~I.~P.~Rubinstein, and C.~Leckie. Labelled data
collection for anomaly detection in wireless sensor networks. In
\textit{2013 6th International Conference on Intelligent Sensors, Sensor
Networks and Information Processing}, pages 269--274. IEEE, 2013.

\bibitem{ref17}
L.~Xie, K.~Lin, S.~Wang, F.~Wang, and J.~Zhou. Differentially private
generative adversarial network. \textit{arXiv preprint arXiv:1802.06739},
2018.

\bibitem{ref18}
J.~Xin, K.~Hu, C.~Luo, M.~Li, Z.~Xie, G.~Mo, Y.~Huang, and C.~Zheng.
Snowtrail: Testing with production traffic. In \textit{Proceedings of the
2022 International Conference on Management of Data}, pages 2272--2284,
2022.

\end{thebibliography}

\end{document}
